% BLOQUE B -- Comparación de herramientas (responsable: Lavayen).
% Datos de la tabla tomados de las páginas oficiales citadas.
% Evidencia propia: Issues #1 (RF-01) y #2 (TC-01) del repositorio del equipo.
\begin{frame}{Herramienta pesada frente a ligera: Helix ALM y GitHub Issues}
  \begin{columns}[T]
    \begin{column}{0.55\textwidth}
      \scriptsize
      \renewcommand{\arraystretch}{1.15}
      \begin{tabular}{@{}p{1.5cm}p{2.6cm}p{2.6cm}@{}}
        \toprule
         & \textbf{Helix ALM}\newline{\tiny(hoy Perforce ALM)} & \textbf{GitHub Issues} \\
        \midrule
        Enfoque & Módulos de requisitos, pruebas e incidencias \cite{perforcealm2026} & Seguimiento de tareas, ideas y errores en el repositorio \cite{githubissues2026} \\
        Trazabilidad & Matriz de trazabilidad automática \cite{perforcealm2026} & Referencias por mención (\#n), sub-issues y dependencias \cite{githubissues2026} \\
        Costo & Precio no publicado; prueba de 30 días \cite{perforcealm2026} & Sin costo en nuestro repositorio público (uso propio) \\
        No conviene si\ldots & El equipo tiene presupuesto limitado y plazos de entrega reducidos, ya que es costoso y complejo de usar & El proyecto exige documentación extensa y trazabilidad formal para auditorías \\
        \bottomrule
      \end{tabular}
    \end{column}
    \begin{column}{0.43\textwidth}
      \centering
      \includegraphics[width=\linewidth,height=0.62\textheight,keepaspectratio]{figuras/issues-rf01-recorte.png}

      \vspace{2pt}
      {\scriptsize Ejemplo propio: RF-01 (\#1) enlazado con su prueba TC-01 (\#2).}
    \end{column}
  \end{columns}

  \vspace{4pt}
  \footnotesize\textbf{Conclusión:} para nuestro proyecto de fin de curso conviene GitHub Issues, ya que es el más óptimo para nosotros como estudiantes de la carrera.
\end{frame}
